\par\noindent\textbf{(1)} 设x=A cos(Ωt−δ)，利用cos³θ=(3cosθ+cos3θ)/4并舍去三次谐波，基频方程的同相/正交分量给f cosδ=A[ω0²−Ω²+3βA²/4]、f sinδ=2ζω0ΩA。平方相加得f²=A²\{[ω0²−Ω²+3βA²/4]²+(2ζω0Ω)²\}，它是非线性幅值隐式方程。
\par\noindent\textbf{(2)} 将共振主频近似置同相刚度为零，得Ω\_res²≈ω0²+3βA²/4。β>0时随A增长频率提高为硬化弹簧，β<0则软化；若幅度太大或失谐导致基频近似差，三次谐波和非周期行为不可忽略。
\par\noindent\textbf{(3)} 令y=A²、d=ω0²−Ω²、b=3β/4、γ=2ζω0Ω，响应方程f²=y[(d+by)²+γ²]。在固定Ω下对y求导，转折满足(d+by)(d+3by)+γ²=0；要有两个实折叠点需d²>3γ²且正根y>0，形成多值响应与跳跃滞回。
\par\noindent\textbf{(4)} 动力学乘ẋ作周期平均，质量和保守非线性项均为全导数，得f⟨cosΩt ẋ⟩=2ζω0⟨ẋ²⟩。单谐波时右侧=ζω0Ω²A²，左侧=fΩA sinδ/2，利用f sinδ=2ζω0ΩA恰相等；该验证保证相位符号与耗散率一致。
